% Propietario conservado: /Users/ruben/Documents/New project/output/REVISION_ARGUMENTAL_APERTURAS_CIERRES_20260915/VII/delivery/MANUSCRIPT_VII_ES_EN_COMPLETE/source_es/manuscrito/30d_realizacion_geometrica.tex
% RESULTADO_RECUPERADO: S15, 11c_gravedad_holonomia_rev6.tex:128--212;
% S01, 02_dinamica_tres_hojas.tex:155--219; pantalla y respuesta de30.
% Pruebas de curvatura, covariancia y Bianchi reunidas para autonomía.
% La realización suave conserva sus condiciones de holonomía explícitas.

\section{Geometric realization of transport and soldering}
\label{vii:vii:sec:realizacion-cartan}

The discrete representation of return and screen soldering
determine two aspects of transport: the accumulation of memory along
routes and the comparison of their fibres. Their geometric realization
preserves both operations. This section fixes the target
space, the dimensional normalization of the connection and the identities
that its variation will use.

\subsection{Screen, co-tetrad and connection}

The screen quotient \(Q_\square\), its Lorentzian form, and the
soldering \(\beta_m\) come from
\ref{vii:vii:sec:pantalla-respuesta}. In each cell, a screen realization
\(J_{{\rm scr},m}\) transports the form and soldering to the
corresponding geometric fibre. Its cellular image is
\[
 e_m=J_{{\rm scr},m}\beta_m.
\]
The scale \(\ell_m\) remains within the soldering. Changes
of frame act jointly on \(J_{{\rm scr},m}\), the connection
and the transport of the cells.

A smooth realization of those data on an oriented manifold
\(\mathcal U\) of dimension four consists of a nondegenerate
co-tetrad \(e=(e^I)\), a Lorentz connection
\(\omega=(\omega^I{}_J)\) and a realization of routes
\(\mathfrak R_{\rm C}\) satisfying
\begin{equation}
 \operatorname{Hol}_{\mathcal A_{\ell_0}}
       (\mathfrak R_{\rm C}(\gamma))
 =\rho_{\rm C}(\operatorname{Hol}_{\rm TPK}(\gamma)).
 \label{vii:vii:eq:compatibilidad-holonomia-geometrica}
\end{equation}
In this matrix equality, the representation is normalized by
\[
 N_{\ell_0}(R,a)=(\Lambda(R),a/\ell_0),\qquad
 \rho_{\rm C}=N_{\ell_0}\circ\rho_{\rm C}^{\rm disc},
\]
where \(\Lambda:\operatorname{Spin}^+(3,1)\to SO^+(3,1)\)
is the vector representation. The map \(N_{\ell_0}\)
preserves the semidirect product: dividing the translation
\(a+\Lambda(R)b\) by \(\ell_0\) is equivalent to composing the
two normalized translations. The complete turn therefore publishes
\(72u\) in the normalized connection and \(72\ell_0u\) in the
dimensional chart. The spin lift separately preserves
\(R_4^{c_g(\gamma)}\); the vector representation has the
central kernel \(\{1,-1\}\), so that lift is retained
when using spinor fields.
Identity, concatenation, inversion and subdivision are maintained in the
corresponding realization. The screen pairing
and the holonomy condition specify the data of the
realization used in the field equations.

With \(\eta=\operatorname{diag}(-1,1,1,1)\), define
\(g=\eta_{IJ}e^I\otimes e^J\). For a positive elementary scale
\(\ell_0\), constant in the chart, the affine connection is
\begin{equation}
 \mathcal A_{\ell_0}
 =\begin{pmatrix}\omega&\ell_0^{-1}e\\0&0\end{pmatrix}.
 \label{vii:vii:eq:conexion-afin-realizada}
\end{equation}
The co-tetrad has the dimension of length; the factor
\(\ell_0^{-1}\) makes its translational component dimensionally homogeneous.

\begin{proposition}[Curvature and torsion of the realized connection]
\label{vii:vii:prop:curvatura-afin}
The curvature of \eqref{vii:vii:eq:conexion-afin-realizada} is
\begin{equation}
 \mathcal F_{\ell_0}
 =\mathrm d\mathcal A_{\ell_0}
   +\mathcal A_{\ell_0}\wedge\mathcal A_{\ell_0}
 =\begin{pmatrix}F_\omega&\ell_0^{-1}T\\0&0\end{pmatrix},
 \quad F_\omega=\mathrm d\omega+\omega\wedge\omega,
 \quad T=\mathrm de+\omega\wedge e.
 \label{vii:vii:eq:curvatura-afin-realizada}
\end{equation}
\end{proposition}
\begin{proof}
The product of matrices of forms produces
\(\omega\wedge\omega\) in the diagonal block and
\(\ell_0^{-1}\omega\wedge e\) in the translational block.
The exterior derivative contributes
\(\mathrm d\omega\) and \(\ell_0^{-1}\mathrm de\), since
\(\ell_0\) is constant. Their sum proves the equality.
\end{proof}

Curvature describes the frame-transport defect and torsion
that of the co-tetrad. Both belong to the same realized connection;
their components retain distinct tensorial domains.

\begin{proposition}[Covariance and Bianchi identities]
\label{vii:vii:prop:bianchi-cartan}
Under a change of frame \(g_L:\mathcal U\to SO^+(3,1)\),
\[
 e'=g_Le,\qquad
 \omega'=g_L\omega g_L^{-1}-\mathrm dg_L\,g_L^{-1},
\]
we have \(T'=g_LT\) and \(F_{\omega'}=g_LF_\omega g_L^{-1}\).
Furthermore,
\begin{equation}
 D_\omega T=F_\omega\wedge e,\qquad D_\omega F_\omega=0.
 \label{vii:vii:eq:bianchi-realizacion}
\end{equation}
\end{proposition}
\begin{proof}
In \(\mathrm d(g_Le)+\omega'\wedge g_Le\), the terms
containing \(\mathrm dg_L\wedge e\) cancel.
The same substitution into the curvature gives its transformation by
conjugation. For the first Bianchi identity, expand
\[
 D_\omega(\mathrm de+\omega\wedge e)
 =(\mathrm d\omega+\omega\wedge\omega)\wedge e.
\]
In \(\mathrm dF_\omega+\omega\wedge F_\omega-F_\omega\wedge\omega\),
the terms containing \(\mathrm d\omega\) and the cubic
products cancel in pairs. This proves the second identity.
\end{proof}

\subsection{The three quadratic realizations of curvature}

The TRIT sheet \(\tau\in\{+1,0,-1\}\) is represented in
the Cartan algebra by generators \(J_{ab}=-J_{ba}\)
and \(P_a^{(\tau)}\), with relations
\begin{align*}
 [J_{ab},J_{cd}]&=
 \eta_{bc}J_{ad}-\eta_{ac}J_{bd}
 -\eta_{bd}J_{ac}+\eta_{ad}J_{bc},\\
 [J_{ab},P_c^{(\tau)}]&=
 \eta_{bc}P_a^{(\tau)}-\eta_{ac}P_b^{(\tau)},\\
 [P_a^{(\tau)},P_b^{(\tau)}]&=-\tau J_{ab}.
\end{align*}
For constant \(\ell>0\), write
\[
 \mathcal A_\tau^{\rm Cart}
 =\tfrac12\omega^{ab}J_{ab}+\ell^{-1}e^aP_a^{(\tau)}.
\]

\begin{proposition}[Curvature of the tritified realization]
\label{vii:vii:prop:curvatura-tritificada}
With the preceding relations,
\begin{equation}
 \mathcal F_\tau^{\rm Cart}
 =\tfrac12(F_\omega^{ab}-\tau\ell^{-2}e^a\wedge e^b)J_{ab}
     +\ell^{-1}T^aP_a^{(\tau)}.
 \label{vii:vii:eq:curvatura-tritificada}
\end{equation}
\end{proposition}
\begin{proof}
Expand \(\mathrm d\mathcal A+\tfrac12[\mathcal A,\mathcal A]\).
The bracket of the Lorentz sector produces \(F_\omega\), the
mixed bracket produces \(D_\omega e\), and the translational
bracket contributes
\(-\tfrac12\tau\ell^{-2}e^a\wedge e^bJ_{ab}\).
\end{proof}

On the central sheet the affine connection is recovered; the extreme sheets
preserve the two signs of the constant-curvature term. The
representation thus retains the tritic regime of the route before
applying the variational equations. The holonomy equality
\eqref{vii:vii:eq:compatibilidad-holonomia-geometrica} corresponds to the
central sheet. On each extreme sheet, the representation
\(\rho_{{\rm C},\tau}\) in its Cartan group and its respective holonomy
condition are used. The current and the metric tensor use
the fields and representation of the same realization.
